\section{Evidence Bundle and Local Packaging}

ASIEP separates the evidence record from the artifacts it references. This separation is necessary because raw traces, feedback comments, candidate diffs, score reports, gate reports, prompts, outputs, and system configurations may be large, sensitive, or access-controlled. Embedding all of them directly in the evidence object would make the record difficult to share, harder to govern, and less reusable across research settings. The bundle design keeps the ASIEP record compact while binding each external artifact through a role, media type, relative path, and digest. The evidence object remains portable without pretending that every source artifact is safe to expose.

A local evidence bundle contains a bundle manifest, an ASIEP evidence record, and artifact files. The manifest declares expected artifacts and records their roles, locations, media types, and digests. The resolver checks whether every referenced artifact is declared, located inside the bundle root, and digest-matched. It also detects missing artifacts, undeclared references, unused artifacts, digest mismatches, and path escape attempts. The resolver is intentionally local-only. It does not fetch remote URIs, authenticate with external services, or resolve globally registered identifiers. This boundary keeps the evaluation reproducible and avoids overstating the implementation as deployed evidence infrastructure.

Packaging is performed only after ASIEP validation and bundle resolution succeed. The local packager then produces four outputs. First, it generates a package manifest that records package identity, included artifacts, digests, and packaging metadata. Second, it creates an FDO-like local record that represents the evidence package as a digital object without claiming registry registration or PID issuance. Third, it generates RO-Crate-like metadata to describe the evidence bundle as a machine-actionable research object. Fourth, it creates PROV JSON-LD to express provenance relations among the improvement cycle, base blueprint, candidate blueprint, evidence artifacts, evaluator activity, and gate report. These outputs make the package easier for another agent or reviewer to inspect, but they do not constitute external standard certification.

The packaging layer is important for eScience-oriented reproducibility. A reviewer can rerun the importer, resolver, validator, packager, and evaluator on the local corpus, inspect generated packages, and compare the resulting metrics with the evaluation report. The package is therefore a reproducible local research object for the current implementation. It supports FAIR-oriented evidence reuse by making artifacts findable within the bundle, accessible under declared local policies, interoperable through mappings to PROV/FDO-like/RO-Crate-like structures, and reusable through explicit roles, digests, and lifecycle metadata.

This design keeps ASIEP's claims modest. The evidence bundle is not a production archive, not a full compliance record, and not a substitute for external governance review. The FDO-like and RO-Crate-like outputs are local representations used to demonstrate packaging feasibility. The PROV JSON-LD output provides a provenance-oriented view of the improvement event, but some mappings are necessarily lossy because ASIEP includes self-improvement-specific constraints, such as gate thresholds and rollback references, that are not native to generic provenance models. The intended result is a locally verifiable evidence package, not a certified infrastructure deployment.
